\section*{1. Ausgangslage / Stand der Forschung}
Die Legacy-Code-Base und der dazugehörige Testpool des Entwicklungsteams sind über die Jahrzehnte gewachsen, veraltet und durch eine starke Verteilung der Umgebung unübersichtlich. Die Ergebnisse des Testpools müssen von verschiedenen Stellen zusammengetragen werden, wodurch das tägliche Monitoring viel Zeit beansprucht.
Um bestehende Produkte zukunftsfähig zu machen, sind Modernisierungen stets gefragt. In dieser Umgebung muss dafür die Übersichtlichkeit gesteigert und der Arbeitsaufwand verringert werden.

Moderne Standards sehen eine übersichtliche, kompakte und nutzerzentrierte Lösung vor. Dabei müssen wichtige Kenngrößen präsent dargestellt sein und auf einem Blick erfasst werden können. Ein konsistentes Design, welches u.a. verwandte Metriken gruppiert, ist ebenso sinnvoll.

Ein neuer Standard trifft auf viele verschiedene Meinungen. Da das Team mit seinem Produkt mitgewachsen ist, kamen bereits vor Projektstart der erste Unmut und Kritik auf. Es wurde die Sorge geäußert, ob die bisherige Lösung nicht mehr ausreicht und diese, trotz Funktionstüchtigkeit, vollständig ersetzt werden soll. Um die Akzeptanz zu steigern und den Übergang zu erleichtern, ist eine nutzernahe Entwicklung entscheidend.

Hierfür wird mithilfe einer kurzen Feedbackschleife gearbeitet, welche in Form einer Nutzerbefragung umgesetzt wird. Diese wird von einem Fragebogen zur Evaluation der Software-Usability von Prof. Dr. Jochen Prümper abgeleitet. Der Fragebogen wurde entwickelt, um die sieben Gestaltungsgrundsätze der DIN EN ISO 9241-110 sicherstellen zu können.
